\begin{tcolorbox}[title= \name{}s Prompt]
\small
\textcolor{teal}{[CLAIM]} \\
Prime Minister Narendra Modi breached the election protocol by addressing a rally in Howrah on April 6.
\textcolor{teal}{[END]}

\textcolor{orange}{[PASSAGE]}\\
Modi addressed a rally in Cooch Behar and Howrah’s Dumurjola on April 6, where polls were held on April 10. The silence period was not breached. ...... In Dumurjola, the voting took place in the fourth phase on April 10. Therefore, Modi did not breach the 48 hours silence period. The voting for West Bengal assembly elections for 294 seats is taking place from March 27 to April 29. \dots\\
Published: 2021-04-07. 
\textcolor{orange}{[END]} 

\textcolor{blue}{[QUERIES]}
1. Modi rally in Howrah 2021

     2. Prime Minister Narendra Modi rally on April 6
     
      3. election in Howrah 2021
      
       4. locations of Howrah voting
        
         5. silence period affects campaigning and media coverage of elections
\textcolor{blue}{[END]}

Using the above as an example, generate at most 10 independent, short, and diverse Google search phrases required to verify or debunk the below claim labeled under \textcolor{teal}{[CLAIM]}.
Note: Generate diverse questions tending to the different numerical and temporal aspects both implicit and explicit to the claim.
\\ Note: Use the passage under \textcolor{orange}{[PASSAGE]} for reference to generate queries. 
%\\ Do not generate queries for the passage.  
%\\ Note: You must exclude any questions about fact-checking. 
\end{tcolorbox}
\captionof{figure}{ Prompt used to generate sub-queries given claim and justification document in the data creation pipeline of \quantempplus{} \\}
\label{example:prompt_agq} 
\section{Data Construction Approach}

An overview of our data collection pipeline  is shown in Figure \ref{fig:data-collection}. We primarily focus on curating a realistic evidence collection without temporal leakage or ``gold" justification leakage.  Additionally to cover all aspects of the claim we propose a claim decomposition process that generates queries around different aspects of the claims by extracting search patterns from justification document of manual fact-checkers. We observe that these queries help in collecting relevant evidence published prior to the date the claim was made. Compared to prior evidence collection approaches \cite{programFC,claimdecomp,multifc}, it also helps reduce noisy evidence stemming from surface level matches to original claim and provides better coverage for different aspects of the claim.

\subsection{\name{}: Automatic Query Generation Emulating Fact-checker Search Process}
\label{fcdecomp}


Manual fact checkers often search multiple news sources or documents on the open web to fetch relevant evidence articles to verify claims. They perform this search by usually generating search queries around different aspects of the claim to ensure the collected evidences are sufficient to verify the complete claim. This search process though not explicitly mentioned on all fact-checking websites is implicit in the justification provided by fact-checkers for verdicts assigned to previously fact-checked claims.

To emulate their thought process, we use few-shot learning to generate these web search queries using the justification document provided for each claim. We hypothesize that utilizing a claim's justification document in a few-shot learning setup with a Large Language Model (LLM) would generate both explicit and implicit queries required to verify a natural numerical claim. 
% LLM’s, specifically Open AI’s GPT models have shown to be effective to simulate user search queries\cite{9999111,10.1145/3366423.3380193}, hence we use the gpt-3.5-turbo model to generate the questions.
The prompt to generate Web search queries includes the claim, its publication date, justification document, and one static example crafted manually by the authors. The instruction specifies the generator to produce queries that sufficiently address the \textbf{numerical aspects} of the claim. The full prompt can be found in the Figure.\ref{example:prompt_agq}. 

To mitigate hallucination and redundant queries, we formulate a filtering process that ensures only queries which are unique and relevant are retained. The filter is constructed using the Maximal Marginal Relevance (MMR) measure.
% Within the algorithm, we employ the paraphrase-MiniLM-L6-v2 model to assess the similarity between texts. 
Given that the queries generated may be explicit or implicit we need to ensure that the selected queries are still relevant but don’t deviate from the claim's core aspects. Therefore, we employ a weighted average approach, to combine the MMR scores with respect to the claim and justification document to form the final scores for each query. Finally, we apply a threshold to retain only the most relevant and independent queries for each claim. Hence, for a given claim \( C \), justification document \( J \) and their corresponding sub-queries \( q_1, \ldots, q_k \), the final set of queries \( Q \) is formed by:
\[
s_{q_i} = \alpha \cdot \text{MMR}(q_{i}, C,\lambda) + (1-\alpha) \cdot \text{MMR}(q_{i}, J,\lambda) ; \quad Q = \{ q_i \mid s_{q_i} > \gamma \text{ for } i \in \{1, k\} \}
\]

% s

Here, \(\alpha\)  is the coefficient used to control the importance needed to be given to relevance with respect to the claim versus that to the justification document, \(\lambda\) is the diversification factor used in MMR and \(\gamma\) is the threshold for selecting a query.

\name{} outputs a variable number of diverse, high-quality queries emulating a manual fact-checker's search process described in the \textbf{justification document} to relevant evidence per claim. The average number of queries per claim is 6.76, which is higher than approaches like PgmFC\cite{programFC} and ClaimDecomp\cite{claimdecomp}. 

\subsection{Collecting Evidence}

After the unique high quality queries are retained, we collect search results from public search engines. However, unlike prior works \cite{multifc,schlichtkrull2023averitec,claimdecomp} which do not account for evidence leakage from fact-checker websites, we filter out evidence from over 150 fact-checking domains adopting the list from \cite{quantemp}. This ensures that the justification written by fact-checkers for the claims does not leak into evidence collected to prevent fact-checking models from learning shortcuts.  While QuanTemp \cite{quantemp} also focuses on preventing such leakage, it does not account for \textit{temporal leakage} where  evidence snippets published after the date the claim was made. This does not reflect realistic fact-checking scenario as systems and manual fact-checkers would only have access to relevant evidence prior to when the claim is made.  To prevent temporal leakage, we add the \textbf{before: filter} to the query issued to the search engine APIs, which restricts the results to those published before the date provided. Then we also perform basic de-duplication to avoid persisting redundant evidences in the collection. The evidence snippets collected from web using these queries are mapped to the claim as relevant evidences after basic de-duplication and also retaining documents that have high diversity as measured by MMR, to form qrels. Since a realistic open-domain retrieval setup also comprises distractors in the corpus, we retain other noisy evidence snippets retrieved from the web prior to filtering to form qrels in the corpus


\subsection{Distilling \name{} Queries and Evidence Aggregation}

While, \name{} helps generate high quality queries, to collect evidence, the queries cannot be used at test time. This is primarily because they rely on justification document for fact-checkers which is not available at test-time in a  realistic automated fact-checking setup. Hence, we propose to train a Language Model to generate such queries by using the claim and corresponding queries generated for the training set using a LLM.

Given the original Claim $C$ from training set, we employ a LLM like \gptt{} to generate a structured list of sub-queries 
$Q = \{q_1, q_2, \ldots, q_n\}$ as described in Section \ref{fcdecomp}.

Since, fine-tuning and inference using another LLM for query generation would be expensive and infeasible, we train a relatively smaller language model $f_\theta$ (e.g., FLAN-T5-LARGE : 780M parameters) 
to imitate the LLM's decomposition behavior. 
The model takes as input the claim $c$ (optionally prefixed with an instruction, 
e.g., ``\textit{Decompose the following claim into search queries:}'') 
and is trained to output the corresponding list of queries $Q$.
The model is optimized using the standard next-token prediction (cross-entropy) loss:
\[
\mathcal{L}(\theta) = - \sum_{(c, Q) \in \mathcal{D}_{\text{train}}} 
\sum_{t=1}^{|Q|} \log P_\theta \big(y_t \mid y_{<t}, c \big),
\]
where $y_t$ is the $t$-th token in the target query list $Q$. At inference time, the model generates a list of sub-queries given the instruction and the claim. 
\vspace{-1em}
\subsection{ Evidence Aggregation for Claim Verification}
The existing claim decomposition approaches like Claimdecomp \cite{claimdecomp}, ProgramFC \cite{programFC} primarily generate queries or sub-claims from original claim followed by retrieval of top-k evidences per query/sub-claim. Then simple concatenation is employed to merge different retrieval lists for claim verification without considering relative ordering. Rank fusion approaches can help in a principled combination of evidences from queries/sub-claims \cite{irfusion,dark_horse_error,irfusion2}.  We evaluate different such approaches and compare them with heuristic top-1 and concatenation based methods for all baselines and \name{}. We observe that while CombSUM and CombMAX provide improvements over naive top-k or concatenation approaches, CombMAX-Norm which applies CombMAX to normalized relevance scores works best to merge evidences across queries from decomposed claims.  %Finally, for claim-verification we employ fine-tuned NLI models or few-shot prompted LLMs given the claim and aggregated evidence snippets.